\documentclass[11pt]{article}

\usepackage[margin=1in]{geometry}
\usepackage{amsmath,amssymb,amsthm,mathtools,booktabs,array,tabularx}
\usepackage{microtype,xurl}
\usepackage{xcolor}
\usepackage[colorlinks=true,linkcolor=blue!55!black,
            citecolor=blue!55!black,urlcolor=blue!55!black]{hyperref}

\newtheorem{theorem}{Theorem}[section]
\newtheorem{proposition}[theorem]{Proposition}
\newtheorem{lemma}[theorem]{Lemma}
\newtheorem{corollary}[theorem]{Corollary}
\theoremstyle{definition}
\newtheorem{definition}[theorem]{Definition}
\theoremstyle{remark}
\newtheorem{remark}[theorem]{Remark}

\newcommand{\R}{\mathbb R}
\newcommand{\cC}{\mathcal C}
\newcommand{\cL}{\mathcal L}
\newcommand{\id}{\mathrm{id}}
\newcommand{\supp}{\operatorname{supp}}
\newcommand{\Cyc}{\operatorname{Cyc}}
\newcommand{\claimid}[1]{\texttt{#1}}

\title{Relational Clocks and Causal Cartan Contraction through Arity Three\\
in Pure-Weyl Gravity on a Berger Universe}
\author{GPT-5.6.sol\thanks{OpenAI model.  Contributed the research
programme, mathematical direction, derivations, machine verification,
referee analyses, and manuscript.  The project was commissioned and directed
by Asger Alstrup Palm (\texttt{asger@area9.dk}), who initiated the questions,
served as the non-technical orchestrator and corresponding human contact, but
claims no technical contribution.}}
\date{Working draft, July 2026}

\begin{document}
\maketitle

\begin{abstract}
We study whether cylinder-time translation can remain presymplectically null
after a healthy
matter clock and the first nonlinear vertices are included in pure-Weyl
gravity.  On $\R\times S^3$ with a Berger spatial metric we exhibit an open
family of exact, smooth, non-conformally-flat solutions sourced by two
standard-sign conformal scalars.  Their rotating phase is an everywhere
timelike clock, the quartic potential is positive, and the internal clock
momentum is nonzero.  Nevertheless, at fixed couplings the linearized lapse
constraint and compact spatial averaging force its variation to vanish on
every smooth solution tangent.  The covariant identity
$\Omega_{\rm total}(\delta,\cL_D)=\omega\,\delta Q_R$ therefore makes the
cylinder generator $D$ a presymplectic degeneracy on this declared linearized
phase space.  This proves infinitesimal nullity, not integrability of a
nonlinear quotient.

At one exact rational member we then formulate the same conclusion
homologically.  The complete
gauge-fixed classical BV presentation has $54$ rows and a support-local
cyclic retract onto a $26$-row retained complex.  Advanced and retarded chain
contractions lift to all rows.  Differentiating the covariant master action
gives complete arbitrary-input four-dimensional operations $q_2$ and $q_3$,
with the required $L_\infty$ identities, cyclicity and local $D$-equivariance.
The unary and binary Cartan primitives extend to a cyclic ternary primitive:
the arity-three source closes by the graded Jacobi identity, and a causal raw
primitive becomes cyclic under the exact $C_4$ Reynolds projector.  Thus the
$D$ action is causally null-homotopic through arity three.  An independent
reduced-action calculation crosses the producer boundary in eight lapse and
radial/Weyl fourth derivatives, agreeing with sixteen ordered coefficients of
the frozen $q_3$ payload; it is not a second full derivation.

The result is classical: its linear nullity statement holds on the displayed
open fixed-coupling Berger branch, while its all-row BV Taylor statement is
proved at one exact rational representative.  It is not an all-orders
nonlinear theorem, a Hadamard construction, a quantum-master-equation result,
an anomaly calculation, or a claim that a nonzero clock momentum is a freely
variable boundary charge.
\end{abstract}

\section*{Authorship and computational inputs}
The named model author produced the mathematical programme, derivations,
verification code, referee responses, and manuscript.  Asger Alstrup Palm
commissioned and directed the work, initiated the questions, and served as
non-technical orchestrator and corresponding human contact; he claims no
technical contribution.  Repository scripts generate the certificates and
sparse coefficient tables used as computational inputs.  They are not
additional authors.  This article has no generated \TeX{} input.
This attribution describes the research archive.  Any submission to a venue
with a different authorship policy requires a separate accountability review;
that review is not represented by the present working draft.

\tableofcontents

\section{Question, phase space and theorem boundary}\label{sec:intro}

There are two superficially conflicting requirements for a relational clock.
Its phase must evolve and possess a nonzero conjugate momentum, yet the total
Hamiltonian generator should remain a constraint on a closed universe.  A
matter phase with zero momentum is not a clock, while a freely variable
nonzero Hamiltonian charge would make time translation a physical symmetry
rather than proper gauge.  The distinction is made by the pullback of the
charge variation to the allowed solution space, not by the background value
of the charge.  We use the BV description of gauge systems
\cite{BV,HT}, the covariant phase-space pairing \cite{CovariantPhaseSpace},
and causal chain contractions in the sense of Green-hyperbolic complexes
\cite{GreenComplexes}.

We work on
\begin{equation}\label{eq:manifold}
 M=\R\times S^3
\end{equation}
with no spatial boundary.  Couplings are fixed.  Write
$\mathcal E_{\alpha_B,\lambda}(\Phi)=0$ for the covariant equations and
$\bar\Phi$ for the Berger solution.  The charge theorem is evaluated on the
pre-reduced smooth tangent space
\begin{equation}\label{eq:tangent-space}
 \mathcal Z_{\bar\Phi}^{\alpha_B,\lambda}
 :=\ker D\mathcal E_{\alpha_B,\lambda}\big|_{\bar\Phi}
 \subset\Gamma^\infty(M,F),
 \qquad \delta\alpha_B=\delta\lambda=0,
\end{equation}
including the linearized constraints.  We test whether the $D$ vector lies in
the radical of the covariant presymplectic form on this space; we do not
quotient by $D$ in the definition.  The nonlinear statement is instead a
formal Taylor identity for the classical BV vector field through arity three
at the same background.  These are related statements, but they are not
identical.

Let $D:=\partial_t$ denote cylinder-time diffeomorphism and let $R$ generate
the global, ungauged internal $O(2)$ rotation of the scalar doublet.  On the
background,
\begin{equation}\label{eq:D-R-definition}
 \cL_D\bar g=0,\qquad \cL_D\bar T=\omega R\bar T,
 \qquad K:=D-\omega R,qquad \cL_K\bar\Phi=0.
\end{equation}
Thus the solution is stationary in stress energy but not pointwise fixed by
$D$ on the scalar fields.  The orbit of $D$ is helical in spacetime times the
internal target.  The paper tests the diffeomorphism generator $D$, not the
stabilizer $K$ and not the global rotation $R$.

\begin{theorem}[Invariant scoped Berger theorem]\label{thm:main}
Fix $\alpha_B>0$.  For every point of the positive Berger branch described in
Section~\ref{sec:background}, assertions (i)--(ii) below hold.  At the exact
rational representative \eqref{eq:rational-fixture}, assertions (iii)--(iv)
also hold:
\begin{enumerate}
\item the background carries a smooth standard-sign conformal-scalar phase
clock with positive quartic potential, timelike phase gradient and nonzero
internal momentum;
\item on the smooth fixed-coupling linearized solution tangent
$\mathcal Z_{\bar\Phi}^{\alpha_B,\lambda}$,
\begin{equation}\label{eq:main-degeneracy}
 \Omega_{\rm total}(\delta,\cL_D)=0
\end{equation}
for every allowed tangent $\delta$;
\item the complete $54$-row gauge-fixed classical BV complex admits
advanced and retarded chain contractions with causal support;
\item if $q_1,q_2,q_3$ are the first three action-derived Taylor operations,
then there are cyclic causal Cartan components
$\iota_D^{(1)},\iota_D^{(2)},\iota_D^{(3)}$ satisfying the Cartan equation
through arity three; equivalently,
\begin{equation}\label{eq:invariant-cartan}
 \bigl([Q,\iota_D]-L_D\bigr)^{(n)}=0,
 \qquad n=1,2,3,
\end{equation}
with an uncomputed arity-four-and-higher remainder.
\end{enumerate}
Consequently $D$ is presymplectically null at the linearized covariant level
and is causally null-homotopic through the first ternary BV interaction test.
No constant-rank neighbourhood, integrated null foliation, finite nonlinear
quotient or all-orders moment-map theorem is asserted.
\end{theorem}

The four assertions are certified respectively by
\claimid{P09-C1}--\claimid{P09-C2},
\claimid{P09-C3}--\claimid{P09-C4},
\claimid{P09-C5}--\claimid{P09-C6}, and
\claimid{P09-C7}--\claimid{P09-C10} in the companion computational
supplement.  The theorem does not assert that the Cartan contraction extends
to all Taylor orders or to a renormalized quantum theory.
The manuscript status is therefore \texttt{WRITING\_STARTED}, not
\texttt{THEOREM\_FROZEN}.

\paragraph{Proof architecture.}
The geometric layer derives the exact background and charge identity from the
action.  The analytic layer constructs the retained Green homotopies and
lifts them through a support-local deformation retract.  The algebraic layer
differentiates the master action and solves the first three Cartan recurrences.
Computer algebra certifies the instantiated curved coefficients and all-row
signs; it does not replace the invariant definitions or the compact-averaging
argument.

\paragraph{Novelty and relation to established machinery.}
Relational and complete observables, covariant phase-space charges,
linearization stability, BV--BFV theory, cyclic homotopy algebras, homotopy
transfer, Cartan homotopies and Green homotopies are established frameworks
\cite{Rovelli,Dittrich,DynamicalFrames,Khavkine,AMM,CMR,Kajiura,
HomotopyTransfer,FiorenzaManetti,GreenComplexes}.  The contribution claimed
here is their explicit conjunction in one action-derived example: a positive
fully backreacting Berger clock, a fixed-coupling presymplectic-nullity
theorem, a complete gauge-fixed classical BV presentation, and a cyclic
causally supported Cartan primitive through $q_3$.  To our knowledge the
specific Berger/Weyl/scalar construction is new; no priority claim is made
for any of the general frameworks.

\section{A positive clock on an exact Berger background}\label{sec:background}

Let $\sigma_i$ be a left-invariant coframe on $S^3$ normalized by
\begin{equation}\label{eq:MC}
 d\sigma_1=-\sigma_2\wedge\sigma_3,qquad
 d\sigma_2=-\sigma_3\wedge\sigma_1,qquad
 d\sigma_3=-\sigma_1\wedge\sigma_2,qquad
 \int_{S^3}\sigma_1\wedge\sigma_2\wedge\sigma_3=16\pi^2.
\end{equation}
We use metric signature $(-,+,+,+)$.  The curvature convention is fixed
operationally by
\begin{equation}\label{eq:curvature-convention}
 \delta\!\int\sqrt{-g}\,C^2
 =4\int\sqrt{-g}\,B_{\mu\nu}\,\delta g^{\mu\nu},
\end{equation}
which, with the matter stress convention below, gives
$\alpha_BB_{\mu\nu}=T_{\mu\nu}$.  Set
\begin{equation}\label{eq:berger}
 g=-dt^2+a^2(\sigma_1^2+\sigma_2^2)+c^2\sigma_3^2,
 \qquad q=\frac{c^2}{a^2}.
\end{equation}
The matter fields are two real conformal scalars with positive target metric,
written in polar form as
\begin{equation}\label{eq:clock}
 T_1=\rho\cos(\omega t),\qquad
 T_2=\rho\sin(\omega t).
\end{equation}
Their phase $\theta=\omega t\pmod{2\pi}$ obeys
$g^{-1}(d\theta,d\theta)=-\omega^2<0$ and is therefore an everywhere
timelike $S^1$-valued clock candidate.  The temporal-diffeomorphism/Weyl incidence determinant
is $\rho^2\omega$, so the clock slice is regular on the branch below.

Our action convention is
\begin{equation}\label{eq:action}
 S=\int_M\!\sqrt{-g}\left[
 \frac{\alpha_B}{8}C_{\mu\nu\rho\sigma}C^{\mu\nu\rho\sigma}
 -\frac12\sum_{A=1}^2\nabla_\mu T_A\nabla^\mu T_A
 -\frac{R}{12}\rho^2-\frac{\lambda}{4}\rho^4
 \right].
\end{equation}
With $T_{\mu\nu}=-2\delta S_m/\delta g^{\mu\nu}$, its metric equation is
\begin{equation}\label{eq:bach-source}
 \alpha_B B_{\mu\nu}=T_{\mu\nu}.
\end{equation}

\begin{proposition}[Exact positive Berger family; \claimid{P09-C1},
\claimid{P09-C2}]\label{prop:background}
For
\begin{equation}\label{eq:q-range}
 \frac{5-\sqrt{21}}2<q<\frac14
\end{equation}
there is an exact solution of the scalar and metric equations with
\begin{equation}\label{eq:branch}
 \rho^2=\frac{2\alpha_B(1-4q)}{a^2},\qquad
 \omega^2=\frac{q}{4a^2(1-4q)},\qquad
 \lambda=-\frac{q^2-5q+1}{6\alpha_B(1-4q)^2}>0.
\end{equation}
The scalar kinetic metric is positive, the quartic potential is bounded
below, and the exact stress obeys the pointwise dominant-energy inequalities
$0\leq p_h,p_v\leq\varepsilon$ in the orthonormal frame.
\end{proposition}

\begin{proof}
In the orthonormal Berger frame the independent Bach components are
\begin{align}
 B_{00}&=\frac{(a^2-c^2)^2}{6a^8},\nonumber\\
 B_{11}=B_{22}&=\frac{(a^2-c^2)(a^2-3c^2)}{6a^8},\label{eq:bach-components}\\
 B_{33}&=\frac{(a^2-c^2)(5c^2-a^2)}{6a^8}.
\end{align}
Substitution of \eqref{eq:clock} into the scalar equation and the three
independent components of \eqref{eq:bach-source} gives
\eqref{eq:branch}.  The interval \eqref{eq:q-range} makes $\rho^2$,
$\omega^2$ and $\lambda$ positive.  Direct substitution gives
\begin{align}
 \varepsilon&=\frac{\rho^2(1-q)^2}{12a^2(1-4q)},\nonumber\\
 p_h&=\frac{\rho^2(1-q)(1-3q)}{12a^2(1-4q)},\qquad
 p_v=\frac{\rho^2(1-q)(5q-1)}{12a^2(1-4q)},
\end{align}
 from which $0\le p_h,p_v\le\varepsilon$ follows on the stated
interval.  The coefficientwise verification also checks the lapse,
horizontal-scale and vertical-scale variations of the Weyl-square action.
\end{proof}

These properties define the word ``healthy'' in this paper: standard-sign
target kinetic form, nonzero amplitude, timelike phase gradient, positive
bounded-below potential on the solution, and the displayed pointwise
dominant-energy inequalities.  They are properties of this exact solution,
not a claim that the improved conformal-scalar stress satisfies an energy
condition on arbitrary configurations \cite{CCJ}.  The inequalities are
pointwise and on shell in the displayed Weyl representative; they are not
asserted to be invariant under changing that representative.

\begin{remark}[$S^1$ clock and local real lifts]\label{rem:S1-clock}
Because $\rho>0$, the map $(T_1,T_2)/\rho$ is globally $S^1$-valued and has
period $2\pi/\omega$ along the $D$ orbit.  A real clock coordinate exists only
after choosing a local lift of $\theta$ on an interval shorter than one
period.  The slice $\theta=\theta_0$ therefore meets a complete orbit
periodically.  We prove transversality of this local clock slice, not a global
one-intersection deparametrization and not absence of Gribov repetitions.
\end{remark}

The rational point
\begin{equation}\label{eq:rational-fixture}
 a=1,\quad q=\frac9{40},\quad \alpha_B=5,\quad
 \rho=1,\quad \omega=\frac34,\quad \lambda=\frac{119}{480}
\end{equation}
is used for the exact all-row BV computations.  It is a representative of
the open branch, not a finite-mode replacement for the four-dimensional
operations.

\section{Nonzero clock momentum and vanishing tangent charge}\label{sec:charge}

The matter action has an internal $O(2)$ rotation
\begin{equation}
 R(T_1,T_2)=(-T_2,T_1)
\end{equation}
with current
\begin{equation}\label{eq:current}
 j_R^\mu=T_1\nabla^\mu T_2-T_2\nabla^\mu T_1.
\end{equation}
For the Maurer--Cartan volume convention
$\int_{S^3}\sigma_1\wedge\sigma_2\wedge\sigma_3=16\pi^2$, the background
charge is
\begin{equation}\label{eq:QR}
 Q_R=16\pi^2\alpha_B q\sqrt{1-4q}>0.
\end{equation}
This is the canonical momentum conjugate to the phase $\theta$.  Thus the
clock evolves and carries a genuine momentum.

The cylinder generator $D=\partial_t$ acts trivially on the stationary metric
and rotates the scalar pair, as already distinguished from $R$ and $K$ in
\eqref{eq:D-R-definition}:
\begin{equation}\label{eq:helical}
 \cL_Dg=0,\qquad \cL_DT=\omega RT.
\end{equation}
The action-derived covariant presymplectic form then gives the exact identity
\begin{equation}\label{eq:omega-charge}
 \Omega_{\rm total}(\delta,\cL_D)=\omega\,\delta Q_R.
\end{equation}
The pure-Weyl cross term vanishes because $\cL_Dg=0$; the matter term is the
variation of the Noether charge \eqref{eq:QR}.  This establishes
\claimid{P09-C3}, but by itself it does not decide whether $D$ is gauge.  The
decision depends on the allowed tangent space.

\subsection{The fixed-coupling slice and the lapse normalization}

To define that tangent space without gauge-fixing away its constraint, retain
the homogeneous lapse and use the common Weyl scale to set $a(t)=1$ only
after variation:
\begin{equation}\label{eq:reduced-metric}
 ds^2=-N(t)^2dt^2+\sigma_1^2+\sigma_2^2+c(t)^2\sigma_3^2,
 \qquad T=\rho(t)(\cos\theta(t),\sin\theta(t)).
\end{equation}
With $V_0=16\pi^2$, the reduced action is
\begin{equation}\label{eq:reduced-action}
 \frac{S_{\rm red}}{V_0}=\int dt\,Nc\left[
 \frac{\alpha_B}{8}C^2
+\frac{\dot\rho^2+\rho^2\dot\theta^2}{2N^2}
-\frac{R\rho^2}{12}-\frac{\lambda\rho^4}{4}\right].
\end{equation}
For auditability, the curvature scalars before the static specialization are
\begin{align}
 R={}&-\frac{c^3N^3-4cN^3+4\dot c\dot N-4N\ddot c}
 {2cN^3},\label{eq:reduced-R}\\
 C^2={}&\frac{4A_-A_+}{3c^2N^6},\qquad
 A_\pm=c^3N^3-cN^3+\dot c\dot N-N\ddot c
 \pm3c\dot cN^2.\label{eq:reduced-C2}
\end{align}
Define the normalized lapse equation
\begin{equation}\label{eq:EN-definition}
 E_N:=\frac1{V_0}\frac{\delta S_{\rm red}}{\delta N}.
\end{equation}
Equations \eqref{eq:reduced-action}--\eqref{eq:EN-definition} fix the
normalization in the linear identity below.

The open background family is a family across theories, not a modulus of one
fixed theory.  Eliminating the background variables gives
\begin{equation}\label{eq:coupling-polynomial}
 F(q;\alpha_B\lambda)
 =q^2-5q+1+6\alpha_B\lambda(1-4q)^2=0,
\end{equation}
and $\partial_qF$ after substituting the branch value of
$\alpha_B\lambda$ is
$3(6q-1)/(4q-1)$, which never vanishes on \eqref{eq:q-range}.  Schematically,
\begin{equation}\label{eq:fixed-slice-diagram}
 \left\{\bar\Phi(q),\alpha_B,\lambda(q)\right\}_{q\in I}
 \longrightarrow \{(\alpha_B,\lambda)\},
 \qquad
 T_{\bar\Phi}\mathcal S_{\alpha_B,\lambda}
 =\ker D\mathcal E_{\bar\Phi}\cap
 \ker d\alpha_B\cap\ker d\lambda .
\end{equation}
Thus $\delta N,\delta c,\delta\rho,\delta\theta$ are varied in deriving the
equations, but an allowed tangent has $\delta\alpha_B=\delta\lambda=0$ and
must satisfy every linearized equation, including $\delta E_N=0$.

\begin{proposition}[Fixed-coupling charge theorem;
\claimid{P09-C4}]\label{prop:charge}
On the smooth fixed-$\alpha_B$, fixed-$\lambda$ linearized solution space
about the positive Berger clock background,
\begin{equation}\label{eq:deltaQzero}
 \delta Q_R=0.
\end{equation}
Hence \eqref{eq:main-degeneracy} holds and $D$ is presymplectically null on
this linearized phase space.
\end{proposition}

\begin{proof}
Direct variation of \eqref{eq:reduced-action}, followed by restriction to the
background branch, gives
\begin{equation}\label{eq:lapse}
 \delta E_N=-\frac{\alpha_Bq^{3/2}}2\frac{\delta Q_R}{Q_R}.
\end{equation}
Here
\begin{equation}\label{eq:relative-Q}
 \frac{\delta Q_R}{Q_R}
 =\frac{\delta c}{\sqrt q}+2\frac{\delta\rho}{\rho}
 +\frac{\delta\dot\theta}{\omega}-\delta N
\end{equation}
at $a=N=1$.  The coefficient in \eqref{eq:lapse} is nonzero throughout
\eqref{eq:q-range}, so every homogeneous allowed tangent has
$\delta Q_R=0$.  At the rational fixture the same row is
\begin{equation}
 \delta E_N=\frac{27\sqrt{10}}{320}\delta N
 -\frac9{16}\delta c-\frac{9\sqrt{10}}{80}\delta\dot\theta
 -\frac{27\sqrt{10}}{160}\delta\rho
 =-\frac{27\sqrt{10}}{320}\frac{\delta Q_R}{Q_R}.
\end{equation}

By Lemma~\ref{lem:averaging}, the compact spatial isometry group preserves the
background and the
linearized equations, while $\delta Q_R$ is an invariant linear functional.
If a smooth inhomogeneous solution tangent had nonzero $\delta Q_R$, averaging
it over the compact group would produce a homogeneous solution tangent with
the same nonzero value, contradicting \eqref{eq:lapse}.  Thus
\eqref{eq:deltaQzero} holds on the complete declared smooth tangent space.
Equation \eqref{eq:omega-charge} gives the result.
\end{proof}

\begin{lemma}[Compact averaging]\label{lem:averaging}
Let $G=SU(2)_L\times U(1)_R$ be the compact spatial isometry group of the
Berger background and let
$\operatorname{Av}_G\delta\Phi=\int_Gg\cdot\delta\Phi\,d\mu_G(g)$.
Then $\operatorname{Av}_G$ is a continuous map in the smooth Fr\'echet
topology, commutes with the fixed-coupling linearized equations, and obeys
\begin{equation}
 \delta Q_R(\operatorname{Av}_G\delta\Phi)
 =\delta Q_R(\delta\Phi).
\end{equation}
\end{lemma}

\begin{proof}
Naturality of the equations and $G$-invariance of $\bar\Phi$ give
$D\mathcal E_{\bar\Phi}(g\cdot\delta\Phi)
=g\cdot D\mathcal E_{\bar\Phi}(\delta\Phi)$; Haar integration therefore
preserves the kernel and all linearized constraints.  The internal charge is
the integral over the closed spatial slice of a $G$-invariant scalar density,
so its differential is a $G$-invariant continuous linear functional.  Haar
normalization then gives the displayed equality.  Residual-isometry zero
modes are retained by the average; none is removed by a spectral projector.
\end{proof}

The compact averaging step is the matter-clock analogue of the familiar
linearization-stability constraints on compact Cauchy surfaces
\cite{LinearizationInstability,AMM}; here the exact lapse identity
\eqref{eq:lapse} supplies the decisive coefficient.

\begin{remark}[What is fixed]
The interval \eqref{eq:q-range} parametrizes theories through the product
$\alpha_B\lambda$; it is not a modulus inside one fixed theory.  At fixed
couplings the stationary branch does not furnish a freely variable $q$ or
$Q_R$.  There is therefore no contradiction between $Q_R>0$ and
$\delta Q_R=0$ on allowed tangents.
\end{remark}

\section{Gauge explanation and relational evolution}\label{sec:relational}

The phase $\theta$ changes along $D$, but a gauge transformation of the total
system need not leave every partial observable fixed.  Relational quantities
compare another field to the clock.  If $X$ is a scalar quantity with
$D X\ne0$, then locally one may evaluate $X$ on the slice
$\theta=\theta_0$ or form combinations invariant under the simultaneous
shift of $X$ and $\theta$.  The background momentum \eqref{eq:QR} makes this
clock slice nondegenerate, while Proposition~\ref{prop:charge} says that the
infinitesimal cylinder motion is presymplectically null in the compact
fixed-coupling tangent space.

This paper proves that $\theta$ is a regular, backreacting clock candidate and
that $D$ is null in the specified linearized and truncated homological senses.
It does \emph{not} construct a complete observable in the sense of
Rovelli--Dittrich \cite{Rovelli,Dittrich}: no global gauge-invariant map
$X\mapsto X|_{\theta=\theta_0}$, uniqueness domain, or complete-observable
Poisson algebra is asserted.  The periodicity in Remark~\ref{rem:S1-clock}
is one reason to keep that distinction explicit.  Relational Maxwell
redshift fixtures in the repository are follow-up applications and are not
dependencies of Theorem~\ref{thm:main}.

This result is sector dependent.  A boundary Hamiltonian, nonzero flux, or a
phase space allowing $\delta Q_R\ne0$ would change the classification of
$D$.  Nothing here promotes the compact Berger verdict to asymptotically
flat, de Sitter, anti-de Sitter, or bounded regions.

\section{The complete causal BV complex}\label{sec:BV}

Let $(\cC_{54},q_1)$ denote the gauge-fixed classical BV complex at the
rational point \eqref{eq:rational-fixture}.  Its row degrees have
multiplicities
\begin{equation}\label{eq:degrees}
 5\text{ in degree }-1,\qquad 22\text{ in degree }0,\qquad
 22\text{ in degree }1,\qquad 5\text{ in degree }2.
\end{equation}
Table~\ref{tab:rows} summarizes the component implementation.  Here $\tau$
is the clock-adapted time ghost, $\sigma$ the Weyl row, and $(R,\Theta)$ the
dressed amplitude and phase rows.
\begin{table}[ht]
\centering
\footnotesize
\begin{tabularx}{\textwidth}{cXrr}
\toprule
Degree & Bundle rows & Minimal & Nonminimal\\
\midrule
$-1$ & spatial diffeomorphism ghosts $c_i$, temporal ghost $\tau$, Weyl ghost
$\sigma$ & $5$ & $0$\\
$0$ & dressed metric $\widehat h_{\mu\nu}$, $R,\Theta$; multipliers $b_a$;
antighost duals $\bar c_a^*$ & $12$ & $10$\\
$1$ & $\widehat h^*_{\mu\nu},R^*,\Theta^*$; multiplier duals $b_a^*$;
antighosts $\bar c_a$ & $12$ & $10$\\
$2$ & ghost/identity antifields $c_i^*,\tau^*,\sigma^*$ & $5$ & $0$\\
\midrule
& Total & $34$ & $20$\\
\bottomrule
\end{tabularx}
\caption{The complete $54$-component gauge-fixed presentation.  The five
nonminimal labels $a$ comprise three spatial-diffeomorphism, one temporal and
one Weyl row.}
\label{tab:rows}
\end{table}

The invariant retained complex has four bundle rows
\begin{equation}\label{eq:retained-bundles}
 \Gamma(T^*S^3)[-1]\longrightarrow
 \Gamma(S^2T^*M)\longrightarrow
 \Gamma(S^2T^*M)[1]\longrightarrow
 \Gamma(T^*S^3)[2],
\end{equation}
of component ranks $3,10,10,3$.  Its unary blocks are the spatial gauge
generator, the retained action Hessian and the negative formal-adjoint gauge
identity.  The remaining eight minimal clock/temporal/Weyl rows and twenty
nonminimal rows form the $28$-component contractible extension.  Thus the
number $54$ is an all-row completeness audit, whereas the statement
$\cC_{54}\simeq\cC_{26}$ is independent of the chosen component ordering.

A gauge-fermion canonical shear preserves the frozen odd Darboux pairing and
gives a support-local cyclic contraction
\begin{equation}\label{eq:SDR}
 \pi_{\rm cl}\iota_{\rm cl}=\id,\qquad
 \id-\iota_{\rm cl}\pi_{\rm cl}
 =q_1S_{\rm cl}+S_{\rm cl}q_1
\end{equation}
onto the retained complex \eqref{eq:retained-bundles}.  This is claim
\claimid{P09-C5}.

For compactly supported sources, the retained complex has advanced and
retarded degree-$(-1)$ chain contractions
\begin{equation}
 \Lambda_{26,\pm}:\Gamma_c(\cC_{26})
 \longrightarrow\Gamma_{J^\pm}(\cC_{26})[-1].
\end{equation}
Lifting them through \eqref{eq:SDR} gives
\begin{equation}\label{eq:lambda54}
 \Lambda_{54,\pm}
 =S_{\rm cl}+\iota_{\rm cl}\Lambda_{26,\pm}\pi_{\rm cl},
\end{equation}
and therefore
\begin{equation}\label{eq:causal-contraction}
 q_1\Lambda_{54,\pm}+\Lambda_{54,\pm}q_1=\id_{54},
 \qquad
 \supp(\Lambda_{54,\pm}f)\subseteq J^\pm(\supp f).
\end{equation}
Indeed,
\begin{align}
 [q_1,\Lambda_{54,\pm}]
 &=[q_1,S_{\rm cl}]
 +\iota_{\rm cl}[q_{1,26},\Lambda_{26,\pm}]\pi_{\rm cl}\nonumber\\
 &=(1-\iota_{\rm cl}\pi_{\rm cl})
 +\iota_{\rm cl}\pi_{\rm cl}=1.
\end{align}
All maps in the retract are differential or pointwise algebraic.  No inverse
Laplacian, inverse curl, transverse--traceless projector, or mode projector
is used.  The complementary-degree adjoint relation and $D$-equivariance
also lift.  These statements comprise \claimid{P09-C6}.

The cyclic higher primitives below instead map compactly supported inputs to
sections supported in the two-sided causal hull $J=J^+\cup J^-$.  They are
neither the causal difference $\Lambda_+-\Lambda_-$ nor separate
advanced/retarded higher operations.

\section{Action-derived interaction tensors}\label{sec:interactions}

For a fluctuation $\phi$ about $\bar\Phi$, expand the classical BV vector
field in the suspended graded-symmetric factorial convention,
\begin{equation}\label{eq:Qexpansion}
 Q_{\bar\Phi}(\phi)
 =q_1(\phi)+\frac1{2!}q_2(\phi,\phi)
 +\frac1{3!}q_3(\phi,\phi,\phi)+O(\phi^4).
\end{equation}
Thus $q_n$ has $n$ inputs, and $q_3$ is the fourth action derivative paired
with one output.  If $\langle-,-\rangle_{\rm BV}$ denotes the degree-one odd
Darboux pairing, the operations are defined invariantly by the polarized
identity
\begin{equation}\label{eq:variational-qn}
 \left\langle x_0,q_n(x_1,\ldots,x_n)\right\rangle_{\rm BV}
 =\left.
 \frac{\partial^{n+1}}
 {\partial t_0\cdots\partial t_n}
 S_{\rm BV}\!\left[\bar\Phi+\sum_{j=0}^nt_jx_j\right]
 \right|_{t=0},
\end{equation}
with the displayed ordering interpreted in the frozen suspended Koszul
convention.  The local clock dressing and gauge-fermion shear are then applied
as BV-canonical transformations to every input, output and dual row.

The proof pipeline is therefore
\begin{equation}\label{eq:proof-pipeline}
 S_{\rm BV}\xrightarrow{\text{mixed variations}}(q_1,q_2,q_3)
 \xrightarrow{\text{Cartan sources}} A_D^{(1)},A_D^{(2)},A_D^{(3)}
 \xrightarrow{\text{Green/Hom contraction}}\iota_D^{(1)},
 \iota_D^{(2)},\iota_D^{(3)}.
\end{equation}
The operations $q_2$ and $q_3$ are obtained by differentiating the covariant
Weyl--clock master action on arbitrary four-dimensional jets and transporting
the result through the same clock and gauge-fermion canonical maps used at
arity one.  They are not fits to residual modes or a harmonic cutoff.

The arity-two and arity-three parts of $Q^2=0$ are
\begin{align}
 [q_1,q_2]&=0,\label{eq:Linf2}\\
 [q_1,q_3]+\frac12[q_2,q_2]&=0.\label{eq:Linf3}
\end{align}
The common action origin supplies cyclicity with respect to the odd Darboux
pairing.  The stationary invariant-frame generator $D=e_0$ acts linearly,
so its higher Taylor components vanish through this order,
\begin{equation}\label{eq:LDhigher}
 L_D^{(2)}=L_D^{(3)}=0,
\end{equation}
and exact coefficientwise verification gives
\begin{equation}\label{eq:Dderivation}
 [D,q_2]=[D,q_3]=0.
\end{equation}
Claims \claimid{P09-C7} and \claimid{P09-C8} bind these statements to the
portable all-row coefficient exports.

The abstract $L_\infty$ relations follow from the classical master equation;
the large PBW calculation verifies that the concrete curved coefficients,
canonical transports and signs instantiate those relations without a dropped
row.  As a cross-producer check, an independently written homogeneous reduced
action gives eight fourth derivatives in the
$(\widehat h_{00},R)$ sector.  They agree exactly with all sixteen ordered
zero-jet coefficients in the published $\widehat h^*_{00}$ and $R^*$ rows.
This probes both the lapse and dressed radial/Weyl directions without
importing the $q_3$ producer.  It is deliberately not described as a complete
second derivation of the payload.

Arity three is the first test involving both $q_3$ and the composite
$[q_2,q_2]$ channel.  It is not the end of the classical expansion: inverse
metrics, curvature and canonical field redefinitions generally generate
nonzero $q_n$ for $n\ge4$ even though the bare scalar potential is quartic.

\section{Causal Cartan contraction through arity three}\label{sec:Cartan}

Write
\begin{equation}\label{eq:iota-expansion}
 \iota_D=\iota_D^{(1)}+\iota_D^{(2)}+\iota_D^{(3)}+\cdots
\end{equation}
for a degree-$(-1)$ Cartan homotopy.  At unary order,
\begin{equation}\label{eq:Cartan1}
 [q_1,\iota_D^{(1)}]=D.
\end{equation}
A one-sided raw solution follows from \eqref{eq:causal-contraction}; averaging
it with its convention-correct cyclic adjoint yields a cyclic unary
primitive with support in the two-sided causal hull.

For clarity, the chain contraction is used on the multilinear Hom complex,
not as an inverse of a coefficient matrix.  If
$\delta=[q_1,-]$ and $A$ is a homogeneous multilinear operator, set
\begin{equation}\label{eq:Hom-contraction}
 \mathfrak H_\pm(A):=\Lambda_{54,\pm}\circ A.
\end{equation}
The graded derivation rule and $[q_1,\Lambda_{54,\pm}]=1$ give
\begin{equation}
 \delta\mathfrak H_\pm(A)+\mathfrak H_\pm\delta(A)=A.
\end{equation}
This identity includes the $q_1$ action on every input slot; no input term is
discarded.  It is the step that turns a $q_1$-closed Cartan source into a raw
causal primitive.

At binary order define
\begin{equation}\label{eq:A2}
 A_D^{(2)}=[q_2,\iota_D^{(1)}].
\end{equation}
Equations \eqref{eq:Linf2} and \eqref{eq:Dderivation} make $A_D^{(2)}$
$q_1$-closed.  A raw causal primitive is cyclicized with
$\Cyc_3=(1+\tau+\tau^2)/3$, giving
\begin{equation}\label{eq:Cartan2}
 [q_1,\iota_D^{(2)}]=-[q_2,\iota_D^{(1)}].
\end{equation}
This is \claimid{P09-C9}.

The ternary source is
\begin{equation}\label{eq:A3}
 A_D^{(3)}=[q_3,\iota_D^{(1)}]+[q_2,\iota_D^{(2)}]-L_D^{(3)}.
\end{equation}

\begin{lemma}[Closure of the ternary source]\label{lem:A3closed}
The complete arbitrary-input source \eqref{eq:A3} obeys
\begin{equation}
 [q_1,A_D^{(3)}]=0.
\end{equation}
\end{lemma}

\begin{proof}
Let $\delta=[q_1,\,\cdot\,]$.  Using \eqref{eq:Linf3},
\eqref{eq:Cartan1} and \eqref{eq:Cartan2}, the only nonzero formal channels
in $\delta A_D^{(3)}$ are
\begin{equation}\label{eq:closure-channels}
 -\frac12[[q_2,q_2],\iota_D^{(1)}]
 +[q_2,[q_2,\iota_D^{(1)}]]-[q_3,D].
\end{equation}
Graded Jacobi gives
\begin{equation}
 [q_2,[q_2,\iota_D^{(1)}]]
 =\frac12[[q_2,q_2],\iota_D^{(1)}],
\end{equation}
and $[D,q_3]=0$ annihilates the last term.  Thus the normalized coefficients
are $-1/2+1/2=0$ and $0$ in the two independent channels.
\end{proof}

Because cyclic coderivations are closed under their graded bracket,
$A_D^{(3)}$ is cyclic.  The causal chain contraction supplies a raw
primitive
\begin{equation}\label{eq:raw3}
 R^{(3)}=-\mathfrak H_+(A_D^{(3)})
 =-\Lambda_{54,+}\circ A_D^{(3)},
 \qquad \delta R^{(3)}=-A_D^{(3)}.
\end{equation}
Pair the output with the odd Darboux form before cyclicizing.  The resulting
four-linear tensor carries the standard Koszul $C_4$ action and the exact
Reynolds projector
\begin{equation}\label{eq:Cyc4}
 \Cyc_4=\frac14(1+\tau+\tau^2+\tau^3).
\end{equation}
Since $\tau^4=1$, $\Cyc_4^2=\Cyc_4$.  Cyclicity of $q_1$ makes $\Cyc_4$
commute with $\delta$, while cyclicity of the source gives
$\Cyc_4A_D^{(3)}=A_D^{(3)}$.  Therefore
\begin{equation}\label{eq:Cartan3}
 \iota_{D,\mathrm{cyc}}^{(3)}=\Cyc_4R^{(3)},
 \qquad
 [q_1,\iota_{D,\mathrm{cyc}}^{(3)}]=-A_D^{(3)}.
\end{equation}
This proves the ternary part of Theorem~\ref{thm:main} and
\claimid{P09-C10}.

The tensorized $C_4$ convention matters.  The actual $54$-row pairing has
$27$ negatively oriented dual slots.  Its duality involution, degree-one
pairing and odd skew orientation are checked on every row.  The $C_4$ group
law is exhausted over all actual degree classes, covering $978{,}736$
admissible degree-zero row quartets with no defect.

\subsection{Support statement}

The Taylor operations and pairing are support-local.  Green contractions and
their cyclic adjoints are causal.  Consequently
\begin{equation}\label{eq:support3}
 \supp\bigl(\iota_{D,\mathrm{cyc}}^{(3)}(f,g,h)\bigr)
 \subseteq J\bigl(\supp f\cup\supp g\cup\supp h\bigr).
\end{equation}
The cyclic primitive is not asserted to be separately retarded or advanced.
The one-sided chain contractions \eqref{eq:causal-contraction} retain that
stronger property; cyclic averaging generally does not.

\section{Interpretation, failures and open directions}\label{sec:limits}

The result resolves one narrow version of the clock question.  A healthy
matter clock can have $Q_R>0$ while the cylinder generator remains
presymplectically null,
because the constraint fixes $Q_R$ within one theory and annihilates its
pullback differential.  The all-row BV calculation shows that this is not
only a homogeneous minisuperspace coincidence: the linearized covariant
phase space, causal complex and first interaction recurrences give compatible
answers.

The following conclusions are not drawn.
\begin{enumerate}
\item No all-orders nonlinear Cartan contraction is claimed.  Arity four and
higher remain separate gates if the complete master action requires them.
\item No constant-rank nonlinear null foliation, integrated gauge quotient,
global deparametrization or complete relational observable is constructed.
\item No Hadamard two-point function, microlocal spectrum condition,
renormalized composite, anomaly cancellation or quantum master equation is
constructed.
\item No positive graviton Hilbert space or interacting probability theorem
follows from the positive clock matter.
\item No noncompact or bounded phase space is covered.  There $D$ may acquire
a boundary charge or flux and cease to be presymplectically null.
\item Relational-redshift fixtures are applications, not part of the main
theorem unless their independent propagation and backreaction gates close.
\end{enumerate}

The correct continuation is therefore not to advertise an unrestricted
solution of the problem of time.  It is to test whether the same Cartan
structure survives higher Taylor order, whether the $S^1$ clock supports
well-defined complete observables on controlled domains, and whether the
conclusion persists under renormalized Ward identities and physical
boundaries.

\section{Reproducibility and claim map}\label{sec:repro}

The companion file
\path{09-relational-clocks-berger-d-cartan-computational-supplement.tex}
contains the focused claim ledger, exact certificate paths, coefficient
inventory, replay commands and freeze checklist.  Its machine-readable source
is
\path{../d_quotient_classical/certificates/PAPER_09_BERGER_CLAIM_TABLE.json}.
The verifier checks strict Draft~2020--12 schema validity, certificate hashes,
required true and false flags, and occurrence of every claim identifier in
the manuscript sources.

The expensive $q_3$ action expansion is not silently rerun by the paper build.
Its content-addressed Tier-2 receipt records the derivation, canonical
transport, row-bounded identity replay and independent portable consumer.
Fast paper checks import that frozen artifact by hash.  A paper release may be
marked \texttt{THEOREM\_FROZEN} only after clean-tree replay and nonlinear-team
sign-off; this working draft is not yet so marked.

\begin{thebibliography}{99}

\bibitem{Bach}
R.~Bach,
\emph{Zur Weylschen Relativit\"atstheorie und der Weylschen Erweiterung des
Kr\"ummungstensorbegriffs},
Math. Z. \textbf{9} (1921), 110--135.

\bibitem{BV}
I.~A. Batalin and G.~A. Vilkovisky,
\emph{Gauge algebra and quantization},
Phys. Lett. B \textbf{102} (1981), 27--31.

\bibitem{HT}
M.~Henneaux and C.~Teitelboim,
\emph{Quantization of Gauge Systems},
Princeton University Press, 1992.

\bibitem{Rovelli}
C.~Rovelli,
\emph{Partial observables},
Phys. Rev. D \textbf{65} (2002), 124013.

\bibitem{Dittrich}
B.~Dittrich,
\emph{Partial and complete observables for Hamiltonian constrained systems},
Gen. Relativ. Gravit. \textbf{39} (2007), 1891--1927.

\bibitem{DynamicalFrames}
C.~Goeller, P.~A. H\"ohn, and J.~Kirklin,
\emph{Diffeomorphism-invariant observables and dynamical frames in gravity:
reconciling bulk locality with general covariance},
arXiv:2206.01193.

\bibitem{CovariantPhaseSpace}
C.~Crnkowic and E.~Witten,
\emph{Covariant description of canonical formalism in geometrical theories},
in \emph{Three Hundred Years of Gravitation}, Cambridge University Press,
1987.

\bibitem{Khavkine}
I.~Khavkine,
\emph{Covariant phase space, constraints, gauge and the Peierls formula},
Int. J. Mod. Phys. A \textbf{29} (2014), 1430009.

\bibitem{LinearizationInstability}
A.~E. Fischer and J.~E. Marsden,
\emph{Linearization stability of the Einstein equations},
Bull. Amer. Math. Soc. \textbf{79} (1973), 997--1003.

\bibitem{AMM}
J.~M. Arms, J.~E. Marsden, and V.~Moncrief,
\emph{The structure of the space of solutions of Einstein's equations II:
several Killing fields and the Einstein--Yang--Mills equations},
Ann. Phys. \textbf{144} (1982), 81--106.

\bibitem{CMR}
A.~S. Cattaneo, P.~Mnev, and N.~Reshetikhin,
\emph{Classical and quantum Lagrangian field theories with boundary},
PoS CORFU2011 (2012), 044.

\bibitem{Kajiura}
H.~Kajiura,
\emph{Noncommutative homotopy algebras associated with open strings},
Rev. Math. Phys. \textbf{19} (2007), 1--99.

\bibitem{HomotopyTransfer}
A.~S. Arvanitakis, O.~Hohm, C.~Hull, and V.~Lekeu,
\emph{Homotopy transfer and effective field theory I: tree level},
arXiv:2007.07942.

\bibitem{FiorenzaManetti}
D.~Fiorenza and M.~Manetti,
\emph{$L_\infty$ algebras, Cartan homotopies and period maps},
arXiv:math/0605297.

\bibitem{CCJ}
C.~G. Callan, S.~Coleman, and R.~Jackiw,
\emph{A new improved energy--momentum tensor},
Ann. Phys. \textbf{59} (1970), 42--73.

\bibitem{GreenComplexes}
M.~Benini, G.~Musante, and A.~Schenkel,
\emph{Green hyperbolic complexes on Lorentzian manifolds},
Commun. Math. Phys. (2023),
\href{https://doi.org/10.1007/s00220-023-04807-5}
{doi:10.1007/s00220-023-04807-5}.

\end{thebibliography}

\end{document}
